\section{An abstract Buchberger procedure}\label{sec:buchberger-procedure}

The Buchberger procedure is formalized as a relation rather than as a recursive function.

\begin{definition}[Buchberger step and the S-polynomial condition]
\label{def:buchberger-step}
A Buchberger step from $G$ to $H$ consists of $p,q\in G$ and a polynomial $h$
such that
\begin{enumerate}[label=(\roman*),leftmargin=2em]
  \item $h$ is a normal form of $\spoly(p,q)$ modulo $G$;
  \item $h\neq0$ and $\LC(h)$ is a unit; and
  \item $H=G\cup\{h\}$.
\end{enumerate}
We say that $G$ satisfies the \emph{S-polynomial condition} if
$\spoly(p,q)\redstar{G}0$ for all $p,q\in G$.  The corresponding Lean
declarations are
\leandoc{Mathlib/RingTheory/MvPolynomial/BuchbergerAlgorithm}{MonomialOrder.BuchbergerStep}{BuchbergerStep}
and
\leandoc{Mathlib/RingTheory/MvPolynomial/BuchbergerAlgorithm}{MonomialOrder.SPolynomialCondition}{SPolynomialCondition}.
The step relation does not prescribe
which pair $(p,q)$ or which normal form $h$ must be chosen.
\end{definition}

Although Definition~\ref{def:buchberger-step} does not explicitly require
$p$ and $q$ to be nonzero and distinct, every genuine Buchberger step has
witnesses with those properties.  If $p=0$, $q=0$, or $p=q$, then
$\spoly(p,q)=0$.  The zero polynomial has no outgoing reduction, so its only
reachable normal form is zero; this contradicts the requirement $h\neq0$ in a
Buchberger step.

\begin{lstlisting}[float=htbp,style=lean,caption={The nondeterministic Buchberger step relation.},label={lst:buchberger-step}]
def BuchbergerStep
    (G H : Set (MvPolynomial σ R)) : Prop :=
  ∃ p ∈ G, ∃ q ∈ G, ∃ h : MvPolynomial σ R,
    Relation.IsNormalFormOf
        (m.ReducesToSet G)
        (m.sPolynomial p q) h ∧
      h ≠ 0 ∧
        IsUnit (m.leadingCoeff h) ∧
          H = insert h G

def SPolynomialCondition
    (G : Set (MvPolynomial σ R)) : Prop :=
  ∀ p ∈ G, ∀ q ∈ G,
    m.sPolynomial p q ⟶*[m, G] 0
\end{lstlisting}

\begin{proposition}[Ideal preservation and strict leading-term growth]
\label{prop:buchberger-growth}
Let $G,H\subseteq R[X_\sigma]$.
\begin{enumerate}[label=(\roman*),leftmargin=2em]
  \item If there is a Buchberger step from $G$ to $H$, then
  \[
    \ideal{H}=\ideal{G}.
  \]
  \item Assume that every nonzero element of $G$ has unit leading
  coefficient.  If there is a Buchberger step from $G$ to $H$, then
  \[
    \ideal{\LT(g)\mid g\in G}
    \subsetneq
    \ideal{\LT(h')\mid h'\in H}.
  \]
\end{enumerate}
\end{proposition}

For part~(i), the adjoined polynomial $h$ is a reduction result of an
S-polynomial of elements of $G$, and therefore $h\in\ideal{G}$.  For
part~(ii), suppose that $\LT(h)$ belonged to the ideal generated by the
leading terms of $G$.  Under the unit-leading-coefficient hypothesis, this
membership produces a one-step reduction of $h$, contradicting the fact that
$h$ is a nonzero normal form modulo $G$.  The subtype
\leandoc{Mathlib/RingTheory/MvPolynomial/BuchbergerAlgorithm}{MonomialOrder.BuchbergerState}{BuchbergerState}
packages precisely this invariant: every nonzero
polynomial in the current set has unit leading coefficient.

\begin{lstlisting}[style=lean,caption={Noetherian termination of the Buchberger step relation.},label={lst:buchberger-termination}]
def BuchbergerState : Type _ :=
  { G : Set (MvPolynomial σ R) //
      ∀ g ∈ G,
        IsUnit (m.leadingCoeff g) ∨ g = 0 }

def BuchbergerState.Step
    (G H : BuchbergerState m) : Prop :=
  m.BuchbergerStep G.1 H.1

theorem BuchbergerState.wellFounded_flip
    [IsNoetherianRing (MvPolynomial σ R)] :
    WellFounded
      (flip (BuchbergerState.Step (R := R) (m := m)))
\end{lstlisting}

The termination measure sends a Buchberger state $G$ to the ideal
$\ideal{\LT(g)\mid g\in G}$.  Because the state contains the
unit-leading-coefficient invariant, every step between Buchberger states
strictly increases this ideal by Proposition~\ref{prop:buchberger-growth}(ii).
Since a Noetherian ring admits no infinite strictly ascending chain of ideals,
there is no infinite forward chain of Buchberger-state steps.  In Lean, this
is encoded as well-foundedness of the flipped step relation, as shown in
Listing~\ref{lst:buchberger-termination}.

\begin{theorem}[Existence and correctness of a terminal Buchberger state]
\label{thm:terminal-correctness}
Let $G\subseteq R[X_\sigma]$ and let $I=\ideal{G}$.  Assume that $R$ has no
zero divisors, $R[X_\sigma]$ is Noetherian, and every nonzero polynomial in
$R[X_\sigma]$ has unit leading coefficient.  Then no infinite sequence of
Buchberger steps starts from $G$.  Moreover, there exist a Buchberger state
$H$ and a finite sequence of Buchberger steps from $G$ to $H$ such that $H$
has no outgoing Buchberger step and the underlying polynomial set of $H$ is a
Gr\"obner basis of $I$.
\end{theorem}

The Lean theorem also takes the corresponding initial-state invariant as an
explicit hypothesis.  The global coefficient hypothesis in
Theorem~\ref{thm:terminal-correctness} implies this invariant.

If $R=K$ is a field, the no-zero-divisor and unit-leading-coefficient
hypotheses in Theorem~\ref{thm:terminal-correctness} are automatic.  If $K$ is
Noetherian and the index type $\sigma$ is finite, Mathlib's form of
Hilbert's basis theorem proves that \texttt{MvPolynomial }\(\sigma\)\texttt{ K} is Noetherian.

\begin{corollary}[Finite terminal Gr\"obner basis over a field]
\label{cor:finite-terminal}
Let $K$ be a Noetherian field, let $\sigma$ be finite, and let
$G\subseteq K[X_\sigma]$ be finite.  If $I=\ideal{G}$, then there exist a
Buchberger state $H$ whose underlying polynomial set is finite and a finite sequence of Buchberger steps from $G$
to $H$ such that $H$ has no outgoing step and its underlying polynomial set is
a Gr\"obner basis of $I$.
\end{corollary}

The finite-output result follows from two preservation theorems in the Lean
development: one Buchberger step preserves finiteness, and every finite chain
of steps preserves finiteness.  These results are
\leandoc{Mathlib/RingTheory/MvPolynomial/BuchbergerAlgorithm}{MonomialOrder.BuchbergerStep.finite}{MonomialOrder.BuchbergerStep.finite}
and
\leandoc{Mathlib/RingTheory/MvPolynomial/BuchbergerAlgorithm}{MonomialOrder.ReflTransGen_BuchbergerStep.finite}{ReflTransGen_BuchbergerStep.finite},
respectively.

\begin{remark}[What is and is not executable]\label{rem:not-executable}
The finiteness results above do not make the construction executable: the
development does not define a \lean{Finset}-based state, a finite queue of
critical pairs, a pair-selection function, or a normal-form function.
\end{remark}
